% ============================================================
%  Chapter 19: Composition Inflation
%  Source: oxford/publications/sources/composition-inflation.tex
% ============================================================

\epigraph{%
  Each additional dimension of state space\\
  multiplies the number of distinguishable trajectories.%
}{}

\section{Distinguishable Trajectories in $d$-Dimensional Partition Space}
\label{sec:ch19-distinguishable}

\begin{definition}[Composition inflation]
  Let a partition hierarchy have branching factor $b$ and be embedded in a
  $d$-dimensional state space.
  The number of \emph{distinguishable trajectories} from the root to depth $n$ is
  \begin{equation}\label{eq:composition-inflation}
    T(n, d) = d\,(d+1)^{n-1}.
  \end{equation}
\end{definition}

\begin{proof}[Derivation]
  At depth 1 there are $d$ choices of sub-partition.
  At each subsequent depth, the number of choices grows by a factor of $(d+1)$:
  the current cell can either continue in the same direction (1 choice) or
  branch to any of the $d$ sub-cells (giving $d+1$ total choices per step).
  Iterating: $T(n, d) = d \cdot (d+1)^{n-1}$.
\end{proof}

For $d = 3$ (three spatial dimensions):
$T(n, 3) = 3 \cdot 4^{n-1}$, growing as $4^n$.
This exponential growth is \emph{composition inflation}: the number of
distinguishable molecular configurations explodes with increasing depth.

\section{Application to Molecular Structure Space}
\label{sec:ch19-molecular}

A molecule of $N$ heavy atoms in a three-dimensional partition space at depth
$n$ contributes $T(n, 3)^N$ distinguishable configurations.
For a drug-like molecule ($N \sim 30$, $n \sim 4$):
$T(4, 3) = 3 \cdot 4^3 = 192$ configurations per atom,
giving $192^{30} \approx 10^{68}$ total configurations --- far exceeding the
number of synthesisable compounds ($\sim 10^{60}$).

Composition inflation explains the combinatorial explosion of chemical space
and places it on a mathematical foundation.

\section{Application to Spectral Complexity}
\label{sec:ch19-spectral}

The number of distinguishable mass spectra of a molecule at partition depth $n$
is $T(n, 3)$ per fragment ion, growing exponentially with fragmentation depth.
This is the theoretical basis for the generative spectral database
(Chapter~\ref{ch:database}): the database generates spectra by walking the
partition tree at depth $n$, not by exhaustive enumeration.
The $O(k)$ query complexity (independent of database size $N$) follows from
the fact that $T(n,d)$ is determined by depth alone.

%%  SOURCE: integrate full derivations from
%%          oxford/publications/sources/composition-inflation.tex

\bigskip
\begin{tcolorbox}[colback=crownblue!5, colframe=crownblue!60!black,
                  title={\textbf{Chapter 19 Summary}}]
\begin{itemize}[noitemsep]
  \item Distinguishable trajectories in a $d$-dimensional partition space
    grow as $T(n,d) = d(d+1)^{n-1}$.
  \item For $d=3$: $T(n,3) = 3\cdot 4^{n-1}$, exponential in depth.
  \item This composition inflation quantifies chemical space size and
    spectral complexity.
  \item The generative database (Chapter~\ref{ch:database}) exploits
    the tree structure to achieve $O(k)$ query complexity.
\end{itemize}
\end{tcolorbox}
